\documentclass[12pt,a4paper]{article}
\usepackage[utf8]{inputenc}
\usepackage[T1]{fontenc}
\usepackage[italian]{babel}
\usepackage{geometry}
\geometry{margin=2.5cm}
\usepackage{setspace}
\setstretch{1.2}
\usepackage{titlesec}
\usepackage{hyperref}

\titleformat{\section}{\bfseries\large}{\thesection.}{1em}{}
\titleformat{\subsection}{\bfseries}{\thesubsection}{1em}{}

\begin{document}

\begin{center}
{\LARGE \textbf{Genere e apprendimento della matematica nella scuola secondaria di primo grado}}\\[6pt]
\end{center}

\vspace{0.5cm}

\noindent
\textbf{Relatore:} Daniel Doz\\
\textbf{Studente:} Giorgia Leoncilli - Università degli Studi di Trieste\\
\textbf{Contatto:} giorgia.leoncilli@studenti.units.it

\vspace{0.7cm}

\section*{Obiettivo}
Il progetto si propone di analizzare come \textbf{stereotipi e differenze di autopercezione legate al genere} influenzino atteggiamenti, motivazione e rendimento in matematica, con l'obiettivo di promuovere una riflessione consapevole e inclusiva sulla disciplina.

\section*{Finalità}
\begin{itemize}
    \item Favorire una \textbf{consapevolezza critica} sugli stereotipi di genere nelle discipline STEM.  
    \item Stimolare \textbf{interesse e motivazione} verso la matematica, in particolare tra le studentesse.  
    \item Offrire \textbf{strumenti di riflessione e discussione} utili a docenti e studenti.  
\end{itemize}

\section*{Descrizione dell'intervento}
L'attività si articola in tre momenti principali:
\begin{enumerate}
    \item \textbf{Somministrazione di un questionario anonimo} rivolto agli studenti della scuola secondaria di primo grado, volto a indagare ansia matematica, autostima, autoefficacia e atteggiamenti verso la disciplina.  
    \item \textbf{Breve seminario interattivo} dedicato a \textbf{Maryam Mirzakhani} e ad altre matematiche di rilievo storico e contemporaneo, con l'obiettivo di valorizzare modelli femminili positivi e stimolare il dialogo su parità di genere e rappresentazioni nella matematica.
    \item \textbf{Somministrazione di un secondo questionario} circa una settimana dopo l'intervento seminariale identico al primo per analizzare la variazione delle risposte.
\end{enumerate}

I dati raccolti saranno trattati in forma anonima secondo la normativa in materia di privacy e utilizzati esclusivamente a fini di ricerca.

\section*{Restituzione}
È prevista una restituzione dei risultati alla scuola, in \textbf{forma aggregata,} e un eventuale \textbf{momento di confronto con i docenti} per discutere le implicazioni educative.

\section*{Durata prevista}
Due incontri in classe della durata totale di circa \textbf{70 minuti}, comprensivo della compilazione del questionario, del seminario e del questionario conclusivo. Il primo incontro di 60 minuti comprende la somministrazione del primo questionario e del seminario, il secondo incontro della durata di 10 minuti prevede la somministrazione del secondo questionario. 

\vfill
\begin{flushright}
\textit{Trieste, Ottobre 2025}\\[6pt]
\textbf{Giorgia Leoncilli}\\
Università degli Studi di Trieste
\end{flushright}

\end{document}
